\section{What expert grammaticality judgments do}
\label{sec:linguistic-judgment-work}

Expert grammaticality judgments can serve several roles. They can screen examples, test diagnostic contrasts, adjudicate corpus candidates, check paraphrases, evaluate minimal pairs, or assess whether a constructed item successfully isolates the intended variable. These are not the same role as estimating how a population would rate the sentence on a Likert scale.

The methodological burden is correspondingly different. An expert-judgment design should say what qualifies the judge, what the judge was asked to evaluate, what information was available, whether judgments were independent, how disagreements were handled, and whether the judgment bears on form, meaning, register, dialect, processing difficulty, or theoretical classification.

Level discipline matters here. A grammaticality claim, an acceptability-distribution claim, a processing claim, and a social-distribution claim are different claims. Treating every expert judgment as participant data blurs those levels instead of protecting them.
